% math4cs-hw4-set-relation-solution.tex

% !TEX program = xelatex
%%%%%%%%%%%%%%%%%%%%
% see http://mirrors.concertpass.com/tex-archive/macros/latex/contrib/tufte-latex/sample-handout.pdf
% for how to use tufte-handout
\documentclass[a4paper, justified]{tufte-handout}

\input{hw-preamble} % feel free to modify this file if you understand LaTeX well
%%%%%%%%%%%%%%%%%%%%
\title{计算机数学 (3-set: 集合的概念与运算; 4-relation: 关系)}
\me{魏恒峰}{hfwei@hnu.edu.cn}{}{}
\date{2026年04月04日 \quad 习题\\2026年05月05日 \quad 答案}
%%%%%%%%%%%%%%%%%%%%
\begin{document}
\maketitle
%%%%%%%%%%%%%%%%%%%%
\noplagiarism % PLEASE DON'T DELETE THIS LINE!
%%%%%%%%%%%%%%%%%%%%
\begin{abstract}
  % \mfigcap{width = 0.85\textwidth}{figs/George-Boole}{George Boole}
  % \begin{center}{\fcolorbox{blue}{yellow!60}{\parbox{0.65\textwidth}{\large
  %   \begin{itemize}
  %     \item
  %   \end{itemize}}}}
  % \end{center}
\end{abstract}
%%%%%%%%%%%%%%%%%%%%
\beginrequired
%%%%%%%%%%%%%%%

%%%%%%%%%%%%%%%
\begin{problem}[相对补与绝对补]
  请证明,
  \[
    A \cap (B \setminus C) = (A \cap B) \setminus (A \cap C).
  \]
\end{problem}

\begin{proof}
  {\bf 证法一:}
  \begin{align*}
    &(A \cap B) \setminus (A \cap C) \\
    =\; & (A \cap B) \cap \overline{A \cap C} \\
    =\; & (A \cap B) \cap (\overline{A} \cup \overline{C}) \\
    =\; & (A \cap B \cap \overline{A}) \cup (A \cap B \cap \overline{C}) \\
    =\; & \emptyset \cup ((A \cap B) \cap \overline{C}) \\
    =\; & (A \cap B) \setminus C
  \end{align*}

  \noindent{\bf 证法二:}

  对任意 $x$,
  \begin{align*}
    x \in (A \cap B) \setminus (A \cap C)
      &\iff x \in A \cap B \land x \notin A \cap C \\
      &\iff x \in A \land x \in B \land \lnot (x \in A \land x \in C) \\
      &\iff x \in A \land x \in B \land (x \notin A \lor x \notin C) \\
      &\iff (x \in A \land x \in B \land x \notin A) \lor (x \in A \land x \in B \land x \notin C) \\
      &\iff \bot \lor (x \in A \land x \in B \land x \notin C) \\
      &\iff x \in A \land x \in B \land x \notin C \\
      &\iff x \in A \land x \in B \setminus C \\
      &\iff x \in A \cap (B \setminus C).
  \end{align*}
  故
  \[
    A \cap (B \setminus C) = (A \cap B) \setminus (A \cap C).
  \]
\end{proof}
%%%%%%%%%%%%%%%

%%%%%%%%%%%%%%%
\newpage
\begin{problem}[对称差]
  请证明,
  \[
    A \cap (B \oplus C) = (A \cap B) \oplus (A \cap C).
  \]
\end{problem}

\begin{proof}
  \begin{align*}
    A \cap (B \oplus C)
      &= A \cap \big((B \setminus C) \cup (C \setminus B)\big) \\
      &= \big(A \cap (B \setminus C)\big) \cup \big(A \cap (C \setminus B)\big) \\
      &= \big((A \cap B) \setminus (A \cap C)\big)
          \cup \big((A \cap C) \setminus (A \cap B)\big) \\
      &= (A \cap B) \oplus (A \cap C).
  \end{align*}
  其中, 第三步使用了上一题的结论。
\end{proof}
%%%%%%%%%%%%%%%

%%%%%%%%%%%%%%%
\newpage
\begin{problem}[广义并、广义交、德摩根律]
  请化简集合 $A$:
  \[
    A = \mathbb{R} \setminus \bigcap_{n \in \mathbb{Z}^{+}} (\mathbb{R} \setminus \set{-n, -n+1, \cdots, 0, \cdots, n-1, n})
  \]
\end{problem}

\begin{solution}
  记
  \[
    X_n \triangleq \set{-n, -n+1, \cdots, 0, \cdots, n-1, n}\text{。}
  \]

  \setcounter{equation}{0}
  \begin{align*}
    A &= \R \setminus \bigcap_{n \in \Z^{+}} (\R \setminus X_n) \\
    &= \R \setminus \Big(\R \setminus \bigcup_{n \in \Z^{+}} X_n \Big) \\
    &= \R \setminus \Big(\R \setminus \Z \Big) \\
    &= \Z
  \end{align*}
\end{solution}
%%%%%%%%%%%%%%%

%%%%%%%%%%%%%%%
\newpage
\begin{problem}[幂集]
  请证明,~\footnote{不, 我有``幂集''恐惧症。}
  \[
    \ps{A} = \ps{B} \iff A = B.
  \]
\end{problem}

\begin{proof}
  \begin{itemize}
    \item 先证
      \[
        \ps{A} = \ps{B} \implies A = B\text{。}
      \]
      假设 $\ps{A} = \ps{B}$。
      \begin{align*}
        &A \in \ps{A} \\
        \implies &A \in \ps{B} \\
        \implies &A \subseteq B
      \end{align*}
      同理可证 $B \subseteq A$。
      因此, $A = B$。
    \item 再证
      \[
        A = B \implies \ps{A} = \ps{B}\text{。}
      \]
      假设 $A = B$。对于任意 $x$,
      \begin{align*}
        &x \in \ps{A} \\
        \implies &x \subseteq A \\
        \implies &x \subseteq B \\
        \implies &x \in \ps{B}
      \end{align*}
      因此, $\ps{A} \subseteq \ps{B}$。
      同理可证 $\ps{B} \subseteq \ps{A}$。
      因此, $\ps{A} = \ps{B}$。
  \end{itemize}
\end{proof}
%%%%%%%%%%%%%%%

%%%%%%%%%%%%%%%
\newpage
\begin{problem}[关系的运算]
  请证明,
  \[
    R[X_1 \setminus X_2] \supseteq R[X_1] \setminus R[X_2]\text{。}
  \]
  请举例说明 $\supseteq$ 不能替换成 $=$。
\end{problem}

\begin{proof}
    任取 $y$,
  \setcounter{equation}{0}
  \begin{align}
    & y \in R[X_{1}] \setminus R[X_{2}] \\[6pt]
    \iff & y \in R[X_{1}] \land y \notin R[X_{2}] \\[6pt]
    \iff & (\exists x_{1} \in X_{1}.\; (x_{1}, y) \in R)
      \land (\forall x_{2} \in X_{2}.\; (x_{2}, y) \notin R)
      \label{eq:2-3} \\[6pt]
    \implies & \exists x \in X_{1} \setminus X_{2}.\; (x, y) \in R. \\[6pt]
      \label{eq:2-4}
    \iff & y \in R[X_{1} \setminus X_{2}]
  \end{align}
  其中, 在第 (\ref{eq:2-4}) 步可取使得 (\ref{eq:2-3})
  中第一个合取子句成立的某个 $x \in X_{1}$。
  根据 (\ref{eq:2-4}) 中第二个合取子句, $x \notin X_{2}$。
  因此, $x \in X_{1} \setminus X_{2}$。

  举例:
  \[
    R = \set{(x_{1}, y), (x_{2}, y)}
    \qquad X_{1} = \set{x_{1}}
    \qquad X_{2} = \set{x_{2}}
  \]
  \[
    R[X_{1} \setminus X_{2}] = R[\set{x_{1}}] = \set{y}
    \qquad R[X_{1}] \setminus R[X_{2}] = \set{y} \setminus \set{y} = \emptyset
  \]
  \[
    R[X_1 \setminus X_2] \supset R[X_1] \setminus R[X_2]
  \]
\end{proof}
%%%%%%%%%%%%%%%

%%%%%%%%%%%%%%%
\newpage
\begin{problem}[关系的运算]
  请证明,
  \[
    (X \cap Y) \circ Z \subseteq (X \circ Z) \cap (Y \circ Z)\text{。}
  \]
  请举例说明, $\subseteq$ 不能换成 $=$。
\end{problem}

\begin{proof}
  任取 $(a, c)$,
  \setcounter{equation}{0}
  \begin{align*}
    & (a, c) \in (X \cap Y) \circ Z \\[6pt]
    \iff & \exists b.\; (a, b) \in Z \land (b, c) \in X \cap Y \\[6pt]
    \iff & \exists b.\; (a, b) \in Z \land (b, c) \in X \land (b, c) \in Y \\[6pt]
    \implies & (\exists b.\; (a, b) \in Z \land (b, c) \in X)
      \land (\exists b.\; (a, b) \in Z \land (b, c) \in Y) \\[6pt]
    \iff & (a, c) \in X \circ Z \land (a, c) \in Y \circ Z \\[6pt]
    \iff & (a, c) \in (X \circ Z) \cap (Y \circ Z)
  \end{align*}

  举例:
  \[
    X = \set{(b_{1}, c)}
    \qquad Y = \set{(b_{2}, c)}
    \qquad Z = \set{(a, b_{1}), (a, b_{2})}
  \]
  \[
    (X \cap Y) \circ Z = \emptyset \circ Z = \emptyset
    \qquad (X \circ Z) \cap (Y \circ Z) = \set{(a, c)} \cap \set{(a, c)} = \set{(a, c)}
  \]
  \[
    (X \cap Y) \circ Z \subset (X \circ Z) \cap (Y \circ Z).
  \]
\end{proof}
%%%%%%%%%%%%%%%

%%%%%%%%%%%%%%%
\newpage
\begin{problem}[关系的性质]
  请证明,
  \[
    R \text{ 是对称且传递的 } \iff R = R^{-1} \circ R
  \]
\end{problem}

\begin{proof}
  \begin{itemize}
    \item 先证
      \[
        R \text{ 是对称且传递的 } \implies R = R^{-1} \circ R\text{。}
      \]
      假设 $R$ 是对称且传递的。

      因为 $R$ 是对称的, 所以
      \[
        R = R^{-1}\text{。}
      \]

      因为 $R$ 是传递的, 所以
      \[
        R \circ R \subseteq R\text{。}
      \]
      故,
      \[
        R^{-1} \circ R = R \circ R \subseteq R\text{。}
      \]

      其次, 任取 $(a, b) \in R$。

      因为 $R$ 是对称的, 所以 $(b, a) \in R$。

      又因为 $R$ 是传递的, 所以 $(b, b) \in R$ 且 $(b, b) \in R^{-1}$。

      因此, $(a, b) \in R^{-1} \circ R$。故, $R \subseteq R^{-1} \circ R$。
    \item 再证
      \[
        R = R^{-1} \circ R \implies R \text{ 是对称且传递的 }\text{。}
      \]
      假设 $R = R^{-1} \circ R$。

      首先,
      \[
        R^{-1} = (R^{-1} \circ R)^{-1} =
          R^{-1} \circ (R^{-1})^{-1} =
          R^{-1} \circ R = R\text{。}
      \]
      所以, $R$ 是对称的。

      其次~\footnote{这里用到了刚刚证明的 $R^{-1} = R$。},
      \[
        R = R^{-1} \circ R = R \circ R\text{。}
      \]
      所以, $R$ 是传递的。
  \end{itemize}
\end{proof}
%%%%%%%%%%%%%%%

%%%%%%%%%%%%%%%
\newpage
\begin{problem}[等价关系]
  一个自反且传递的二元关系 $R \subseteq X \times X$
  称为 $X$ 上的拟序。
  现令 $\preceq\; \subseteq X \times X$ 为拟序。
  如下定义 $X$ 上的关系 $\sim$:
  \[
    x \sim y \iff x \preceq y \land y \preceq x,
  \]
  请证明, $\sim$ 是 $X$ 上的等价关系。
\end{problem}

\begin{solution}
  \begin{enumerate}[(1)]
    \item $\sim$ 是自反的。\\
      对于任意 $x \in X$, 因为 $\preceq$ 是自反的, 所以
      \[
        x \in X \implies x \preceq x \implies x \preceq x \land x \preceq x
        \implies x \sim x\text{。}
      \]
    \item $\sim$ 是对称的。\\
      对于任意 $(x, y) \in \;\sim$,
      \[
        x \sim y \implies x \preceq y \land y \preceq x
        \implies y \preceq x \land x \preceq y \implies y \sim x\text{。}
      \]
    \item $\sim$ 是传递的。\\
      对于任意 $(x, y) \in \;\sim$ 与 $(y, z) \in \;\sim$,
      \begin{align*}
        & x \sim y \land y \sim z \\[6pt]
        \implies & (x \preceq y \land y \preceq x) \land (y \preceq z \land z \preceq y) \\[6pt]
        \implies & (x \preceq y \land y \preceq z) \land (z \preceq y \land y \preceq x) \\[6pt]
        \implies & x \preceq z \land z \preceq x \qquad \\[6pt]
        \implies & x \sim z\text{。}
      \end{align*}
  \end{enumerate}
\end{solution}
%%%%%%%%%%%%%%%

%%%%%%%%%%%%%%%%%%%%
% 如果没有需要订正的题目，可以把这部分删掉

% \begincorrection
%%%%%%%%%%%%%%%%%%%%

%%%%%%%%%%%%%%%%%%%%
% 如果没有反馈，可以把这部分删掉
% \beginfb

% 你可以写 (也可以发邮件)
% \begin{itemize}
%   \item 对课程及教师的建议与意见
%   \item 教材中不理解的内容
%   \item 希望深入了解的内容
%   \item $\cdots$
% \end{itemize}
%%%%%%%%%%%%%%%%%%%%
\end{document}